%% kapa: exam=1 round=32 year=2021 period=1 form=A subject=civil_law subject_ko=민법 seq=21 number=21
%% source: Q-Net 감정평가사 기출문제 (cst00309) artlSeq=5208397 file=extracted/1008/5208397/A형/2021년도 제32회 감정평가사 1차 1교시 A형.pdf
%% conversion: kordoc 4.24.51 -> tools/gen1.py (auto); review=proofread-sonnet; flags=none
\begin{question}{21}
부동산등기에 관한 설명으로 옳지 않은 것은? (다툼이 있으면 판례에 따름)
\begin{choices}
  \item 전부 멸실한 건물의 보존등기를 신축한 건물의 보존등기로 유용하는 것은 허용된다.
  \item 물권에 관한 등기가 원인 없이 말소되었더라도 특별한 사정이 없는 한 그 물권의 효력에는 아무런 영향을 미치지 않는다.
  \item 소유권이전청구권 보전의 가등기가 있더라도 소유권이전등기를 청구할 어떤 법률관계가 있다고 추정되지 않는다.
  \item 가등기권리자가 가등기에 기한 소유권이전의 본등기를 한 경우에는 등기공무원은 그 가등기 후에 한 제3자 명의의 소유권이전등기를 직권으로 말소하여야 한다.
  \item 소유권이전등기가 마쳐지면 그 등기명의자는 제3자는 물론이고 전소유자에 대해서도 적법한 등기원인에 의하여 소유권을 취득한 것으로 추정된다.
\end{choices}
\end{question}
\answer{1}
